\section{真正的麻烦：无聊、低唤醒、动机、主动切换}

上一章说完了「算法观」，这一章回答一个更实际的问题：ADHD 真正卡住你的，到底是哪几件事？

答案是四件——\term{无聊、低唤醒、动机不足、主动任务切换成本过高}。专注力不够反而是最表面的那个。逐一说。

\subsection{无聊是一种认知负担，不是闲得慌}

普通人对「分心」的默认态度是：分心是敌人，要消灭它。但原视频给出的观察正好相反：\term{对 ADHD 来说，无聊才是大问题，分心不是}（BV1q9PmzxEt2）。

硬坐、硬逼自己专注，对神经连接的负担，比「带着分心做事」大得多。你逼自己盯着一页纸，脑子里念头一个接一个往外冒，每压一个念头都在耗电——而那个念头本来可能是回血用的。

所以 ADHD 的首要任务不是「别分心」，是\term{先别无聊}。原视频甚至建议：先别管定向，优先保证「一天能做得久、做得开心」。无聊解决了，一天干十几个小时也不累；定向的调控，是后面的、漫长的、慢慢练的事。

\subsection{低唤醒：人醒着，脑子部分睡着了}

ADHD 经常经历一种很诡异的疲劳：明明没干什么重活，却觉得浑身没力气；午休睡了两个小时，起来反而更累。这不是矫情，也不是懒，原视频把它解释为\term{低唤醒状态}（BV16mXnBxEE1）——有研究显示，ADHD 成人在清醒状态下，大脑会出现类似睡眠的慢波活动。也就是说：\term{你看着是醒着的，但有一部分脑子已经睡着了}。

这个现象带来两个重要推论：

\begin{itemize}
  \item \term{休息对 ADHD 是「唤醒恢复」，不是「静止恢复」}。被动躺平、午睡、硬睡，对普通人是充电，对 ADHD 可能是放电——没有刺激，大脑进入更昏沉的状态。这也是为什么刷几分钟视频、看到个感兴趣的内容，你反而清醒了。
  \item \term{低唤醒时，执行功能下线，但自然语言和审查能力还在}。你没法线性地看一段有难度的内容，但可以「随便瞄一下」，然后冒出几个问题去问 AI——这是 ADHD 专属的低配模式，照样能干活。
\end{itemize}

针对唤醒度，原视频给了一套三档管理法：

\begin{checklist}
  \item \textbf{高唤醒}：做线性任务，趁热打铁；
  \item \textbf{中唤醒}：随手翻找，看有没有什么激发兴趣；
  \item \textbf{低唤醒}：躺下，允许大脑自发放出想法，用手机或 AI 里的内容重新唤醒，而不是硬撑着执行。
\end{checklist}

\key{你觉得自己「没体力」的时候，往往不是真的没体力，是唤醒度太低——缺的不是休息，是刺激。}

\subsection{意志耗能：抑制一个念头，可能比做那件事更费电}

原视频有个非常反直觉的观察（BV1DQAkzbEuK）：\term{压住一个念头去做别的事，消耗的能量可能比做事本身还大}。

他举的例子是拍视频：脑子里一个视频「读条读满了」，很想录，但按老规矩「早上录完，下午做数学」，于是得把录视频的冲动压一个下午——结果数学也没做成，还多睡了两小时。后来改成「随时写、随时录」，抑制的耗能消失，数学和视频两不误。

把这个放大到生活：取个快递本身三分钟体力，但你拖着不去，是因为出门可能撞见同事、而工作还没做完、怕被评判——\term{压住「怕被评判」的念头，比取快递费电得多}。执行功能在这里扮演的角色，与其说是「干活」，不如说是「维持 + 意志」：维持住目标，压住其他念头。

\begin{insight}
所以 ADHD 的调控方向，不是「提高意志力」，而是\term{降低需要意志的场合}。怎么降？原视频给了一个关键原则：\term{降低权重，而不是排序}。别给自己排一个 ABCD 的执行顺序表——排序本身就是高耗能的抑制过程。把一堆东西在内心合并成一件大事（后面「恰当场」会细讲），让大脑觉得「今天我就做这一件事，其他都不重要」，意志耗能立刻降下来。
\end{insight}

\subsection{主动切换贵，被动切换便宜}

这是整个方法库的地基（BV11ucZzxEkx）：

\term{主动任务切换}（我正做着 A，强迫自己去做 B）＝ 抑制当前冲动，贵得离谱；
\term{被动任务切换}（脑子里的念头自己冒出来，跟着它走）＝ 念头自发涌现，便宜。

所以最优策略不是「按计划切换」，而是\term{把一堆事合并成一条宽的抽象轨道，让念头在轨道里自己跳}。外部看你在频繁切换，内部看你一直在一件事上。这一条的实操化，就是后面方法论里的「恰当场」。

\subsection{假性疲劳：缺的是定向能量，不是追随性能量}

最后一种假性疲劳（BV1MWwJzxEh5）：你累的时候，缺的往往是\term{定向能量}（「我接下来必须写那篇东西」），而不是\term{追随性能量}（「随便看看有什么有意思的」）。

很多人拖着不做正事时的循环是这样的：五点多回宿舍，一想「晚上要看一篇文章」，立刻觉得「好累，先睡一会」，于是从五点多睡到七点、八点、九点……一晚上就这么过去了。\term{「必须要做 X」这个念头本身，就在前额叶持续耗电}，你休息得越久，越起不来。

破解思路是\term{动态调节自由度}：想法上高自由度（「今天不一定要写这个，随便偏出去写点别的也算成果」），操作上低自由度（把要写的文件放在桌面最显眼处，记录一直开着）。这样既取消了「必须」的警报，又保证「啪的一下」的想法能最快落到正事上。

\key{ADHD 的问题，大部分不是意志力问题，是能量管理问题。先把无聊和焦虑解决了，再谈定向。}
